\documentclass[11pt]{article}
\usepackage{amsmath, amssymb, amsthm}
\usepackage[margin=1in]{geometry}

\newtheorem{theorem}{Theorem}
\newtheorem{lemma}{Lemma}
\newtheorem{proposition}{Proposition}
\newtheorem{corollary}{Corollary}

\title{Lower bounds for perfect difference sets, infinitely often \\
  \large Erd\H{o}s Problem \#1194 --- expository note}
\author{}
\date{April 30, 2026}

\begin{document}
\maketitle

\begin{abstract}
We give two self-contained derivations of i.o.\ lower bounds on
perfect difference sets. Theorem~\ref{thm:nlogn} reproves the classical
$a_n \gg n \log n$ infinitely often via Sidon density (Erd\H{o}s,
\cite{HaRo66}). Theorem~\ref{thm:n2io} recovers the recent refinement
$a_n \gg n^{2-o(1)}$ infinitely often, due to GPT-5.4 Pro and
streamlined by Bloom in the comment thread of erdosproblems.com/1194
(April 23 2026). The latter argument does not use Sidon density at all;
it uses only the PDS partition identity and a gap lower bound from
the hypothetical pointwise upper bound on $a_n$. We also flag a
typographical error on the live problem page (``diverges'' vs
``converges'').
\end{abstract}

\section{Setup and notation}\label{sec:setup}

A set $A \subset \mathbb{N}$ is a \emph{perfect difference set} (PDS) if
every $n \geq 1$ admits a \emph{unique} representation
\[
   n = a_n - b_n, \qquad a_n, b_n \in A, \quad a_n > b_n.
\]
The notation $a_n$ is fixed by this definition: it is the larger element
of the unique representing pair for $n$. It is \emph{not} the $n$-th
element of $A$. To avoid confusion, we write the increasing enumeration
of $A$ as
\[
   A = \{x_1 < x_2 < x_3 < \cdots\},
\]
and use $a_n$ exclusively for the representing pair.

A PDS is in particular a Sidon set: all pairwise differences are
distinct. We further write $A_x := A \cap [1, x]$, $k(x) := |A_x|$,
and $N(x) := \#\{n \geq 1 : a_n \leq x\}$.

\section{The PDS counting identity}\label{sec:counting}

\begin{lemma}[PDS counting]\label{lem:pds-count}
For every $x \geq 1$,
\[
   N(x) \;=\; \binom{k(x)}{2}.
\]
\end{lemma}

\begin{proof}
Each $n \leq N(x)$ has $a_n \leq x$ and so $b_n < a_n \leq x$, putting both
elements of the unique representing pair in $A_x$. The map
$n \mapsto \{a_n, b_n\}$ is therefore an injection from $\{n : a_n \leq x\}$
into the unordered pairs of $A_x$, so $N(x) \leq \binom{k(x)}{2}$.
Conversely, every unordered pair $\{a, b\} \subset A_x$ with $a > b$ produces
a difference $a - b \in [1, x-1]$, which by uniqueness is the representation
of some unique $n$ with $a_n = a$, $b_n = b$, and $a_n \leq x$. Hence
$\binom{k(x)}{2} \leq N(x)$.
\end{proof}

\begin{corollary}\label{cor:trivial}
$\binom{k(x)}{2} \leq x$ and so $k(x) \leq (1 + o(1)) \sqrt{2x}$.
\end{corollary}

\noindent
This recovers the standard Sidon density bound; the PDS gives \emph{equality},
not only inequality.

\section{Classical $n \log n$ infinitely often}\label{sec:nlogn}

\begin{theorem}[Erd\H{o}s, after \cite{HaRo66}]\label{thm:sidon-io}
For every infinite Sidon set $B \subset \mathbb{N}$,
\[
   \liminf_{x \to \infty} \frac{|B \cap [1,x]|}{\sqrt{x/\log x}}
   \;<\; \infty.
\]
\end{theorem}

\noindent We use Theorem~\ref{thm:sidon-io} as a black box; the proof is
elementary and standard.

\begin{theorem}\label{thm:nlogn}
For any perfect difference set, $a_n \gg n \log n$ for infinitely many $n$.
\end{theorem}

\begin{proof}
Choose $x_j \to \infty$ with $k(x_j) \leq C \sqrt{x_j / \log x_j}$
by Theorem~\ref{thm:sidon-io}. By Lemma~\ref{lem:pds-count},
\[
   N(x_j) \;=\; \binom{k(x_j)}{2}
   \;\leq\; \frac{C^2 x_j}{2 \log x_j}.
\]
Setting $n_j := N(x_j) + 1$ gives $a_{n_j} > x_j$ and
$n_j \leq C^2 x_j / (2 \log x_j) + 1$. Hence
\[
   a_{n_j} \;>\; x_j
   \;\geq\; \frac{2 \log x_j}{C^2}\,(n_j - 1)
   \;\geq\; \frac{2 (1+o(1))}{C^2}\, n_j \log n_j,
\]
using $\log x_j \geq (1+o(1)) \log a_{n_j} \geq (1+o(1)) \log n_j$.
\end{proof}

\section{The refinement to $n^{2-o(1)}$}\label{sec:n2}

The refinement that follows is due to GPT-5.4 Pro and was condensed by
Thomas Bloom in the comment thread on the problem page. It does not use
Sidon density at all. The argument hinges on two structural facts about
PDS that we extract first.

\subsection{Two preparatory lemmas}

For each $k \geq 2$, let
\[
   D_k \;:=\; \{ x_k - x_i \;:\; 1 \leq i < k \}.
\]

\begin{lemma}[PDS partition]\label{lem:partition}
The family $\{D_k\}_{k \geq 2}$ is a partition of $\mathbb{Z}_{> 0}$.
In particular, for every $X \geq 1$,
\[
   X \;=\; \sum_{k \geq 2} |D_k \cap [1, X]|.
\]
\end{lemma}

\begin{proof}
Every $n \geq 1$ has a unique representation $n = a_n - b_n$ with
$a_n, b_n \in A$, $a_n > b_n$. Pick $k$ such that $a_n = x_k$ and let
$i$ be such that $b_n = x_i$; then $i < k$ and $n = x_k - x_i \in D_k$.
Conversely, every element $x_k - x_i \in D_k$ is a positive integer.
Uniqueness in PDS gives that $n$ lies in exactly one $D_k$.
\end{proof}

\begin{lemma}[Gap lower bound from $a_n \ll n^c$]\label{lem:gap}
Suppose $a_n \leq C\, n^c$ for all $n \geq 1$ and some $c > 1$. Then for
every $k \geq 2$,
\[
   x_k - x_{k-1} \;\geq\; (x_k / C)^{1/c}.
\]
Consequently, $x_k \;\gg\; k^{c/(c-1)}$.
\end{lemma}

\begin{proof}
Let $d_k := x_k - x_{k-1}$. The pair $(x_k, x_{k-1})$ realises the
difference $d_k$; by PDS uniqueness, $a_{d_k} = x_k$. The hypothesis
gives $x_k = a_{d_k} \leq C\, d_k^c$, so $d_k \geq (x_k/C)^{1/c}$.
Telescoping,
\[
   x_k \;\geq\; x_k - x_1
       \;=\; \sum_{i=2}^{k} d_i
       \;\geq\; C^{-1/c} \sum_{i=2}^{k} x_i^{1/c}.
\]
Standard induction (or comparison with the ODE
$y' = C^{-1/c} y^{1/c}$) yields $x_k \gg k^{c/(c-1)}$.
\end{proof}

\subsection{Main contradiction}

\begin{theorem}[GPT-5.4 Pro, after Bloom's exposition]\label{thm:n2io}
For any perfect difference set and any $c < 2$, the inequality
$a_n \leq C n^c$ \emph{cannot} hold for all $n$ and any constant $C$.
Equivalently, for every $c < 2$,
\[
   \limsup_{n \to \infty} \frac{a_n}{n^c} \;=\; \infty,
\]
i.e.\ $a_n \gg n^{2-o(1)}$ infinitely often.
\end{theorem}

\begin{proof}
Suppose, for contradiction, that $a_n \leq C\, n^c$ holds for all $n$,
for some constants $C > 0$ and $c \in (1, 2)$.
Lemma~\ref{lem:gap} then gives $x_k \gg k^{c/(c-1)}$.

Fix $X \geq 1$ and split the partition identity from
Lemma~\ref{lem:partition}:
\begin{equation}\label{eq:split}
   X \;=\; \sum_{k \geq 2} |D_k \cap [1, X]|
       \;=\; \underbrace{\sum_{2 \leq k \leq K} |D_k \cap [1, X]|}_{(\mathrm{I})}
         + \underbrace{\sum_{k > K} |D_k \cap [1, X]|}_{(\mathrm{II})},
\end{equation}
where $K$ will be chosen as $K = \lfloor X^{(c-1)/c} \rfloor$.

\medskip
\noindent \textbf{Bound on (I).} Trivially $|D_k| \leq k - 1$, so
\[
   (\mathrm{I}) \;\leq\; \sum_{k \leq K} (k-1)
              \;\leq\; K^2 / 2 \;\ll\; X^{2 - 2/c}.
\]

\medskip
\noindent \textbf{Bound on (II).} For $k > K$ we have
$x_k \gg k^{c/(c-1)} \gg K^{c/(c-1)} \asymp X$, so $x_k \gg X$ and in
particular $x_k > 2X$ for all sufficiently large $X$ (absorb a constant
into the choice of $K$). The set
\[
   \{x_i : x_i \in [x_k - X, x_k)\}
\]
is in bijection with $D_k \cap [1, X]$. Any two consecutive elements
$x_{j-1} < x_j$ in $[x_k - X, x_k)$ satisfy
\[
   x_j - x_{j-1} \;\geq\; (x_j/C)^{1/c}
                 \;\geq\; ((x_k - X)/C)^{1/c}
                 \;\gg\; x_k^{1/c}
\]
by Lemma~\ref{lem:gap} and $x_k - X \geq x_k/2$. Therefore
\[
   |D_k \cap [1, X]| \;\ll\; \frac{X}{x_k^{1/c}}
                       \;\ll\; X \cdot k^{-1/(c-1)},
\]
again using $x_k \gg k^{c/(c-1)}$. Summing,
\[
   (\mathrm{II}) \;\ll\; X \sum_{k > K} k^{-1/(c-1)}.
\]
Since $c < 2$, we have $1/(c-1) > 1$, so the sum converges; tail
estimation gives
\[
   \sum_{k > K} k^{-1/(c-1)} \;\ll\; K^{1 - 1/(c-1)}
                              \;=\; K^{(c-2)/(c-1)}.
\]
Substituting $K = X^{(c-1)/c}$,
\[
   (\mathrm{II}) \;\ll\; X \cdot X^{(c-1)/c \cdot (c-2)/(c-1)}
                  \;=\; X^{1 + (c-2)/c}
                  \;=\; X^{2 - 2/c}.
\]

\medskip
\noindent \textbf{Combining.} From \eqref{eq:split},
$X \ll X^{2 - 2/c}$, i.e.\
$X^{2/c - 1} \ll 1$. Since $c < 2$, the exponent $2/c - 1 > 0$, so the
left-hand side tends to infinity as $X \to \infty$ --- a contradiction.
\end{proof}

\subsection{Logarithmic refinement}

The inequality $X \ll X^{2-2/c}$ is sharp at $c = 2$, and the contradiction
fails there. By tracking the constants more carefully under a hypothesis
$a_n \leq C\, n^2 / f(n)$ (instead of $a_n \leq C\, n^c$), the same scheme
yields the following.

\begin{theorem}[Logarithmic refinement, after Bloom]\label{thm:n2overf}
Let $f \colon \mathbb{N} \to \mathbb{R}_{>0}$ be nondecreasing, slowly
varying (i.e.\ $f(\lambda n)/f(n) \to 1$ as $n \to \infty$ for every fixed
$\lambda > 0$), unbounded, and satisfy
\[
   \sum_{n \geq 2} \frac{1}{n\, f(n)} \;<\; \infty.
\]
Then for any perfect difference set,
\[
   a_n \;\gg\; \frac{n^2}{f(n)} \qquad \text{for infinitely many } n.
\]
In particular, $a_n \gg n^2 / (\log n)^{1 + \varepsilon}$ infinitely often
for every $\varepsilon > 0$, and $a_n \gg (n / \log n)^2$ infinitely often.
\end{theorem}

\begin{proof}
Suppose, for contradiction, that there exist $C > 0$ and $n_0 \in \mathbb{N}$
with
\begin{equation}\label{eq:hyp}
   a_n \;\leq\; C\, n^2 / f(n)  \qquad \text{for all } n \geq n_0.
\end{equation}

\medskip
\noindent \textbf{Step 1 (gap bound).}
By PDS uniqueness, $a_{d_k} = x_k$ where $d_k = x_k - x_{k-1}$.
For $d_k \geq n_0$, applying \eqref{eq:hyp} gives
$x_k \leq C\, d_k^2 / f(d_k)$, equivalently
\begin{equation}\label{eq:gap2}
   d_k \;\geq\; \bigl( x_k\, f(d_k) / C \bigr)^{1/2}.
\end{equation}

\medskip
\noindent \textbf{Step 2 (size of $x_k$).}
We claim there exist $c_1 > 0$ and $k_0 \in \mathbb{N}$ such that
\begin{equation}\label{eq:xkbound}
   x_k \;\geq\; c_1\, k^2\, f(k) \qquad (k \geq k_0).
\end{equation}

Indeed, by \eqref{eq:gap2} and $d_k \leq x_k$, slow variation of $f$ gives,
for $k$ sufficiently large,
$f(d_k) \geq (1/2) f(\sqrt{x_k})$, hence
$d_k \geq c_2 \sqrt{x_k\, f(\sqrt{x_k})}$ with
$c_2 = (2C)^{-1/2}$. Telescoping,
\[
   x_k \;\geq\; \sum_{i=k_0+1}^k d_i
        \;\geq\; c_2 \sum_{i=k_0+1}^k \sqrt{x_i\, f(\sqrt{x_i})}.
\]
This recurrence is monotone in $x_i$; comparing with the differential
equation $dy/dk = c_2 \sqrt{y\, f(\sqrt{y})}$ (separable as
$dy/\sqrt{y} = c_2 \sqrt{f(\sqrt y)}\, dk$), integrating, and using slow
variation $f(\sqrt y) \asymp f(y^{1/4}) \asymp \cdots$ to absorb the
self-dependence into a constant, we obtain
$\sqrt{x_k} \geq (c_2/2)\, k\, \sqrt{f(k)}\, (1 - o(1))$, hence
\eqref{eq:xkbound} with $c_1 = c_2^2/8$.

\medskip
\noindent \textbf{Step 3 (anchor window).}
Fix $X$ large. Let $K = K(X)$ be the largest integer with $x_K \leq 2X$.
By \eqref{eq:xkbound}, $K^2 f(K) \leq 2X / c_1$, so
\begin{equation}\label{eq:Kbound}
   K \;\leq\; \sqrt{2X / (c_1\, f(K))}
     \;\leq\; \sqrt{2X / (c_1\, f(\sqrt{X/f(\sqrt{X})}))}
     \;\sim\; \sqrt{2X / (c_1\, f(\sqrt X))},
\end{equation}
where the last step uses $f$ slowly varying (since $K = \Theta(\sqrt{X/f(\sqrt X)})$
gives $f(K) \asymp f(\sqrt X)$). Crucially, since $f(\sqrt X) \to \infty$,
\begin{equation}\label{eq:Klittleo}
   K^2 \;=\; o(X) \qquad (X \to \infty).
\end{equation}

For $k > K$ we have $x_k > 2X$. The window $W_k := [x_k - X, x_k)$ contains
some subset of $A$. For any two consecutive $x_{j-1} < x_j$ in $W_k$,
\eqref{eq:gap2} gives $d_j \geq \sqrt{x_j f(d_j)/C}$. Now $x_j \geq x_k - X
\geq x_k/2$, so $\sqrt{x_j} \geq \sqrt{x_k}/\sqrt{2}$; slow variation of
$f$ on $[d_j, x_j]$ gives $f(d_j) \geq (1/2) f(\sqrt{x_k})$ for $k$ large.
Therefore
\[
   d_j \;\geq\; (1/\sqrt{4C})\, \sqrt{x_k\, f(\sqrt{x_k})}.
\]
By \eqref{eq:xkbound} and slow variation, $\sqrt{x_k f(\sqrt{x_k})}
\geq c_3\, k\, f(k)$ for $k$ large, with $c_3 > 0$ depending only on $c_1$.
Hence
\begin{equation}\label{eq:wcount}
   |D_k \cap [1, X]| \;=\; |A \cap W_k|
   \;\leq\; \frac{X}{c_3\, k\, f(k)} + 1.
\end{equation}

\medskip
\noindent \textbf{Step 4 (partition sum).}
By Lemma~\ref{lem:partition},
\[
   X \;=\; \sum_k |D_k \cap [1, X]|.
\]
Split at $K$:
\[
   X \;=\; \underbrace{\sum_{2 \leq k \leq K} |D_k \cap [1, X]|}_{\mathrm{(I)}}
       \;+\; \underbrace{\sum_{k > K} |D_k \cap [1, X]|}_{\mathrm{(II)}}.
\]
For (I), trivially $|D_k| \leq k - 1$, so
$(\mathrm{I}) \leq K(K-1)/2 \leq K^2/2 = o(X)$ by \eqref{eq:Klittleo}.

For (II), apply \eqref{eq:wcount}:
\[
   (\mathrm{II}) \;\leq\; \frac{X}{c_3} \sum_{k > K} \frac{1}{k\, f(k)}
                   \;+\; M_K,
\]
where $M_K = |\{k > K : |A \cap W_k| \geq 1\}|$ counts which large-$k$
windows are non-empty. Each non-empty $W_k$ contributes a distinct
$x \in A$ (since $W_k$ ranges over disjoint anchor positions $x_k$),
and any two consecutive $x_k, x_{k+1}$ in this counting differ by
$d_{k+1} \geq c_4 \sqrt{x_k f(\sqrt{x_k})} \geq c_4 \sqrt X$ (using
$x_k > 2X$ and slow variation). So $M_K \leq c_5 \sqrt{X}$ (the number
of $A$-anchors that could possibly have a window intersecting $[1, X]$
is bounded by Sidon density of $A$). In particular $M_K = O(\sqrt X) = o(X)$.

Let $\varepsilon_K := \sum_{k > K} 1/(k f(k))$. By the convergence
hypothesis, $\varepsilon_K \to 0$ as $K \to \infty$, and in particular
$\varepsilon_K \to 0$ as $X \to \infty$. Thus
\[
   (\mathrm{II}) \;\leq\; \frac{X \varepsilon_K}{c_3} + o(X).
\]

\medskip
\noindent \textbf{Step 5 (contradiction).}
Combining,
\[
   X \;\leq\; \frac{K^2}{2} \;+\; \frac{X \varepsilon_K}{c_3} \;+\; o(X)
     \;=\; \frac{X \varepsilon_K}{c_3} \;+\; o(X).
\]
For $X$ large, $\varepsilon_K < c_3/2$ and the $o(X)$ term is less than
$X/4$, so
\[
   X \;\leq\; \frac{X}{2} + \frac{X}{4} \;=\; \frac{3X}{4},
\]
which is false. This contradicts \eqref{eq:hyp}.

Hence \eqref{eq:hyp} fails for every $C$: for each $C > 0$, there are
infinitely many $n$ with $a_n > C\, n^2/f(n)$, i.e.\ $a_n \gg n^2/f(n)$
infinitely often.
\end{proof}

\begin{remark}
The case $f(n) = (\log n)^{1+\varepsilon}$ ($\varepsilon > 0$) gives
$a_n \gg n^2/(\log n)^{1+\varepsilon}$ infinitely often; taking
$\varepsilon \to 0^+$ recovers $a_n \gg n^{2-o(1)}$ from
Theorem~\ref{thm:n2io}. The convergence boundary
$\sum 1/(n f(n)) < \infty$ is sharp \emph{for this argument}: at
$f(n) = \log n$ (where the sum diverges, $\sum 1/(n \log n) = \infty$),
the tail $\varepsilon_K$ in Step~4 no longer goes to zero with $K$ and
the contradiction collapses. Pushing past this boundary would require
either a sharper density estimate on the anchor window in Step~3 or a
non-partition counting identity.
\end{remark}

\begin{remark}[\emph{Typographical note on the live problem page}]
\label{rem:typo}
At time of writing (April 30, 2026), the problem page
\texttt{erdosproblems.com/1194} states the conclusion of
Theorem~\ref{thm:n2overf} as ``for any function $f$ such that
$\sum 1/(n f(n))$ \emph{diverges}.'' This appears to be a copy-edit
error: \emph{convergence} of $\sum 1/(nf(n))$ is the natural sufficient
condition for the proof, and is what Bloom states (with example
$f(n) = (\log n)^2$) in the comment thread after natso26's correction
(post 5782). With the divergence condition, $f(n) = 1$ would qualify and
imply $a_n \gg n^2$ infinitely often, which the proof does \emph{not}
establish.
\end{remark}

\section{What is \emph{not} given by these arguments}\label{sec:caveats}

Theorems~\ref{thm:nlogn} and~\ref{thm:n2io} are \emph{infinitely often}
statements: they exhibit a sparse subsequence $n_j \to \infty$ on which
$a_{n_j}$ is large. They place no constraint on $a_n$ for ``most'' $n$.

Empirically (see \texttt{verify\_and\_analyze.py}), the running minimum
$\inf_{n \leq N} a_n / n^2$ tends to zero in the greedy PDS:
$\min_{n \leq 1898} a_n/n^2 \approx 5.6 \times 10^{-4}$. So a
``for-all-large-$n$'' version of $a_n \gg n^2$ is empirically false even
for the greedy PDS, and the original Erd\H{o}s question is best
reformulated as a statement on the running maximum
\[
   M(N) \;:=\; \max_{n \leq N} a_n.
\]
Theorems~\ref{thm:nlogn} and~\ref{thm:n2io} translate into
$M(N) \gg N \log N$ and $M(N) \gg N^{2-o(1)}$ respectively. Lev's
greedy upper bound is $M(N) = O(N^3)$, refined empirically to
$M(N) \approx 0.0293\, N^3$ for $N \in [1500, 1898]$.

The unresolved question is therefore: \emph{is there a constant
$\alpha \in (2, 3]$ such that $M(N) \asymp N^\alpha$? In particular,
can the upper bound be brought down from $N^3$, or the lower bound
brought up to match?} The Cilleruelo--Nathanson method
\cite{CiNa08} for transferring results from dense Sidon sets to PDS
is a candidate route on the upper side.

\begin{thebibliography}{99}
\bibitem{HaRo66} H.~Halberstam and K.~F.~Roth, \emph{Sequences},
  Oxford University Press, 1966.
\bibitem{Le04} V.~F.~Lev, \emph{Reconstructing integer sets from their
  representation functions}, Electron.\ J.\ Combin.\ 11 (2004), R78.
\bibitem{CiNa08} J.~Cilleruelo and M.~B.~Nathanson, \emph{Perfect
  difference sets constructed from Sidon sets},
  Combinatorica 28 (2008), 401--414.
\end{thebibliography}

\end{document}
